% ============================================================
%  chapters/ch2.tex — Chương 2: Cơ sở lý thuyết
% ============================================================
\section{CƠ SỞ LÝ THUYẾT}
\label{sec:ch2}

% ============================================================
\subsection{Giao thức RPL}
\label{subsec:rpl}
% ============================================================

\subsubsection{Tổng quan về RPL}

Giao thức \textbf{RPL} (\textit{IPv6 Routing Protocol for Low-Power and
Lossy Networks}) được IETF chuẩn hóa theo RFC~6550 năm 2012, là giao thức
định tuyến chuẩn \textit{de facto} cho các mạng IoT năng lượng thấp và mất
mát cao (LLN) \cite{winter2012rpl}. RPL hoạt động ở tầng mạng, theo hướng
tiếp cận định tuyến vector khoảng cách (distance vector), và tổ chức mạng
thành cấu trúc đồ thị acyclic có hướng tập trung vào đích (DODAG).

RPL được thiết kế để đáp ứng bốn mô hình truyền thông chính trong IoT:

\begin{itemize}
    \item \textbf{Nhiều-đến-một (Many-to-One):} Các nút cảm biến gửi dữ
          liệu về nút gốc. Đây là mô hình phổ biến nhất trong thu thập
          dữ liệu cảm biến.
    \item \textbf{Một-đến-nhiều (One-to-Many):} Nút gốc phát lệnh hoặc
          cấu hình xuống các nút cảm biến (downward routing).
    \item \textbf{Điểm-đến-điểm (Point-to-Point):} Hai nút bất kỳ trao
          đổi trực tiếp thông qua nút gốc.
    \item \textbf{Đa điểm-đến-đa điểm (Multipoint-to-Multipoint):} Hỗ
          trợ qua các cơ chế mở rộng.
\end{itemize}

\subsubsection{Cấu trúc DODAG}

\textbf{DODAG} (\textit{Destination Oriented Directed Acyclic Graph}) là
cấu trúc tô-pô cốt lõi mà RPL xây dựng và duy trì. Mỗi DODAG được đặc
trưng bởi:

\begin{itemize}
    \item \textbf{DODAG Root (Nút gốc / Border Router -- BR):} Nút đặc
          biệt kết nối LLN với mạng IP bên ngoài. Mọi luồng upward đều
          hướng về nút này.
    \item \textbf{DODAGID:} Địa chỉ IPv6 toàn cục duy nhất định danh
          một DODAG.
    \item \textbf{DODAGVersionNumber:} Bộ đếm 8-bit, tăng mỗi khi nút
          gốc khởi tạo tái cấu trúc tô-pô toàn cục (global repair).
    \item \textbf{Rank:} Giá trị số nguyên không âm biểu diễn vị trí
          tương đối của nút so với nút gốc. Nút gốc có Rank = 1; các
          nút càng xa gốc có Rank càng lớn. Rank đóng vai trò quan trọng
          trong việc ngăn chặn vòng lặp định tuyến.
\end{itemize}

Mỗi nút không phải gốc phải chọn một \textbf{nút cha ưu tiên}
(\textit{preferred parent}) trong số các nút có Rank nhỏ hơn. Tập hợp
các nút cha khả dĩ tạo thành \textit{parent set}. Luồng dữ liệu upward
đi từ nút lá (leaf node) qua các nút cha lên tới nút gốc theo cấu trúc
cây phân cấp.

\subsubsection{Các bản tin điều khiển RPL}

RPL sử dụng bốn loại bản tin điều khiển ICMPv6 để xây dựng và duy trì
DODAG \cite{prajisha2025dis}:

\paragraph{DODAG Information Object (DIO)}
DIO là bản tin quảng bá định tuyến quan trọng nhất. Nút gốc phát DIO
multicast định kỳ; các nút trung gian sau khi nhận DIO sẽ cập nhật thông
tin và chuyển tiếp DIO xuống các nút con. DIO chứa các trường chính:
DODAGID, DODAGVersionNumber, Rank của người gửi, Mode of Operation (MOP)
và thông tin Objective Function.

\paragraph{DODAG Information Solicitation (DIS)}
DIS là bản tin yêu cầu thông tin DODAG, được một nút gửi khi muốn gia
nhập hoặc tái gia nhập mạng. Khi nhận DIS, các nút hàng xóm đang là thành
viên DODAG sẽ phản hồi bằng DIO và \textbf{reset bộ hẹn giờ Trickle},
dẫn đến phát DIO ngay lập tức thay vì chờ hết chu kỳ hiện tại.

\paragraph{Destination Advertisement Object (DAO)}
DAO được gửi theo hướng upward để thông báo khả năng tiếp cận
(\textit{reachability}) của các nút xuôi dòng. Nút gốc dùng thông tin
DAO để xây dựng bảng định tuyến downward, cho phép gửi dữ liệu từ gốc
xuống bất kỳ nút nào trong DODAG.

\paragraph{DAO Acknowledgment (DAO-ACK)}
DAO-ACK xác nhận việc tiếp nhận và xử lý thành công một bản tin DAO,
đảm bảo độ tin cậy trong việc thiết lập đường định tuyến downward.

\subsubsection{Bộ hẹn giờ Trickle}

\textbf{Trickle Timer} là cơ chế điều tiết tần suất phát bản tin DIO,
cho phép RPL thích ứng với điều kiện mạng \cite{prajisha2025dis}. Hoạt
động theo nguyên tắc:

\begin{enumerate}
    \item Mỗi chu kỳ $I$ bắt đầu với khoảng thời gian $I_{min}$ và tăng
          gấp đôi sau mỗi chu kỳ ổn định, tới giới hạn $I_{max}$.
    \item Trong mỗi chu kỳ, nút đếm số lần nghe thấy DIO nhất quán
          (counter $c$). Nếu $c < k$ (redundancy constant), nút phát
          DIO tại thời điểm ngẫu nhiên trong $[I/2, I]$.
    \item Khi phát hiện không nhất quán (nhận DIS, thay đổi parent,
          routing loop...), Trickle Timer bị \textbf{reset} về $I_{min}$,
          kích hoạt phát DIO ngay lập tức để sửa chữa nhanh.
\end{enumerate}

Cơ chế này giảm thiểu overhead điều khiển khi mạng ổn định, nhưng cũng
là điểm yếu bị khai thác bởi DIS Flooding Attack.

\subsubsection{Objective Function}

\textbf{Objective Function (OF)} là tập quy tắc xác định cách tính Rank
từ các routing metrics và cách chọn preferred parent \cite{decrank2022}.
RPL chuẩn hóa hai OF:

\begin{itemize}
    \item \textbf{OF0} (RFC~6552): Tối thiểu hóa số hop đến nút gốc.
          Rank được tính bằng cách cộng một giá trị không đổi
          (rank-increase) vào Rank của preferred parent. OF0 luôn chọn
          nút có Rank thấp nhất làm parent mà không có cơ chế chống
          dao động (anti-churn), khiến nó dễ bị ảnh hưởng bởi
          Decreased Rank Attack.
    \item \textbf{MRHOF} (RFC~6719): Sử dụng metric ETX (Expected
          Transmission Count) để chọn đường có chất lượng link tốt nhất.
          MRHOF tích hợp \textit{hysteresis} --- nút chỉ thay đổi parent
          khi ETX mới tốt hơn ngưỡng $\Delta_{ETX}$ so với parent hiện
          tại, giúp giảm parent churn nhưng vẫn có thể bị tấn công Rank
          nếu $\Delta_{ETX}$ đủ lớn.
\end{itemize}

% ============================================================
\subsection{Phân loại tấn công RPL}
\label{subsec:attack_taxonomy}
% ============================================================

Theo khung phân loại được trình bày trong \cite{overview_rpl_framework},
các tấn công vào RPL được tổ chức thành ba nhóm chính, phản ánh ba chiều
tác động khác nhau lên mạng:

\subsubsection{Nhóm tấn công tài nguyên (Resources)}

Mục tiêu là làm cạn kiệt năng lượng, bộ nhớ hoặc băng thông của các nút
hợp lệ, từ đó rút ngắn tuổi thọ mạng. Nhóm này chia thành hai loại:

\begin{itemize}
    \item \textbf{Tấn công trực tiếp (Direct):} Nút độc hại trực tiếp
          tạo ra lưu lượng quá mức. Ví dụ điển hình: \textit{DIS Flooding
          Attack} --- phát tán DIS liên tục khiến mạng bùng nổ DIO.
    \item \textbf{Tấn công gián tiếp (Indirect):} Nút độc hại kích thích
          các nút hợp lệ tạo ra lưu lượng dư thừa. Ví dụ:
          \textit{Version Number Attack} --- tăng DODAGVersionNumber bất
          hợp pháp buộc toàn mạng tái cấu trúc.
\end{itemize}

\subsubsection{Nhóm tấn công tô-pô (Topology)}

Mục tiêu là phá vỡ cấu trúc DODAG, gây định tuyến sai hoặc cô lập nút.
Chia thành:

\begin{itemize}
    \item \textbf{Sub-optimization:} Mạng hội tụ về trạng thái không
          tối ưu, giảm hiệu năng nhưng không cắt đứt kết nối hoàn toàn.
          Ví dụ: \textit{Decreased Rank Attack}, \textit{Routing Table
          Falsification}.
    \item \textbf{Isolation:} Một hoặc nhiều nút bị cô lập khỏi DODAG.
          Ví dụ: \textit{Blackhole Attack}, \textit{Sinkhole Attack},
          \textit{Wormhole Attack}.
\end{itemize}

\subsubsection{Nhóm tấn công lưu lượng (Traffic)}

Mục tiêu là khai thác hoặc giả mạo dữ liệu lưu thông trong mạng. Chia
thành:

\begin{itemize}
    \item \textbf{Eavesdropping (Passive):} Nghe trộm lưu lượng mà không
          can thiệp. Ví dụ: \textit{Sniffing}, \textit{Traffic Analysis}.
    \item \textbf{Misappropriation (Active):} Chiếm đoạt và thay đổi
          lưu lượng. Ví dụ: \textit{Decreased Rank Attack},
          \textit{Identity Attack}.
\end{itemize}

% ============================================================
\subsection{Tổng quan về IDS trong IoT/LLN}
\label{subsec:ids_overview}
% ============================================================

\subsubsection{Khái niệm và phân loại IDS}

\textbf{Hệ thống Phát hiện Xâm nhập} (IDS) là tập hợp các công cụ và
phương pháp giám sát hoạt động mạng hoặc hệ thống, phát hiện các hành
vi bất thường hoặc vi phạm chính sách bảo mật. Trong môi trường
IoT/LLN, IDS phải đối mặt với thách thức đặc biệt: tài nguyên tính toán
cực kỳ hạn chế, môi trường truyền thông không đáng tin cậy và tô-pô mạng
thay đổi động.

Các IDS cho RPL/LLN được phân loại theo hai chiều chính:

\paragraph{Theo phương pháp phát hiện:}
\begin{itemize}
    \item \textbf{Signature-based:} So sánh hành vi quan sát được với
          cơ sở dữ liệu mẫu tấn công đã biết. Độ chính xác cao, FPR thấp
          nhưng không phát hiện được tấn công mới (zero-day).
    \item \textbf{Anomaly-based:} Xây dựng mô hình hành vi bình thường
          và phát hiện lệch chuẩn. Phát hiện được tấn công mới nhưng
          FPR thường cao hơn.
    \item \textbf{Threshold-based:} Dạng đặc biệt của anomaly-based,
          sử dụng ngưỡng cố định hoặc động cho các metrics cụ thể (tần
          suất DIS, PDR, Rank...). Nhẹ nhàng, phù hợp với LLN.
    \item \textbf{Hybrid:} Kết hợp nhiều phương pháp để tận dụng ưu
          điểm của từng loại \cite{prajisha2025dis}.
\end{itemize}

\paragraph{Theo kiến trúc triển khai:}
\begin{itemize}
    \item \textbf{Centralized:} IDS tập trung tại nút gốc hoặc server
          ngoài. Xử lý mạnh nhưng tạo điểm đơn hỏng (single point of
          failure) và overhead truyền thông cao.
    \item \textbf{Distributed:} Mỗi nút chạy module IDS cục bộ, chia
          sẻ thông tin với hàng xóm. Khả năng mở rộng tốt, phù hợp LLN.
    \item \textbf{Hybrid:} Kết hợp giám sát cục bộ tại mỗi nút với
          phân tích tổng hợp tại nút gốc \cite{prajisha2025dis}.
\end{itemize}

\subsubsection{Thách thức IDS trong LLN}

Thiết kế IDS cho RPL/LLN phải giải quyết các thách thức đặc thù:

\begin{enumerate}
    \item \textbf{Ràng buộc tài nguyên:} Nút cảm biến thường chỉ có
          vài KB RAM, vài chục KB Flash và nguồn pin giới hạn. IDS không
          được tiêu tốn quá nhiều CPU cycles hay bộ nhớ.
    \item \textbf{Môi trường truyền thông không tin cậy:} Packet loss
          tự nhiên cao (LLN) có thể bị nhầm là tấn công Blackhole, gây
          FPR cao.
    \item \textbf{Tô-pô động:} Thay đổi parent hợp lệ (do di chuyển,
          pin yếu) có thể giống với dấu hiệu tấn công.
    \item \textbf{Phân biệt tấn công và lỗi mạng:} Cần cơ chế phân
          biệt hành vi độc hại với lỗi giao thức thông thường.
\end{enumerate}

% ============================================================
\subsection{Các chỉ số đánh giá hiệu năng}
\label{subsec:metrics}
% ============================================================

Để đánh giá toàn diện tác động của tấn công và hiệu quả của IDS, đồ án
sử dụng hai nhóm chỉ số sau:

\subsubsection{Chỉ số hiệu năng mạng}

\paragraph{Packet Delivery Ratio (PDR)}
PDR đo tỉ lệ gói dữ liệu được gửi thành công đến đích trên tổng số gói
được phát đi:
\begin{equation}
    \text{PDR} = \frac{N_{\text{received}}}{N_{\text{sent}}} \times 100\%
    \label{eq:pdr}
\end{equation}
PDR là chỉ số phản ánh trực tiếp nhất mức độ ảnh hưởng của tấn công lên
dịch vụ dữ liệu. Tấn công Blackhole và Decreased Rank làm giảm PDR đáng
kể; tấn công DIS Flooding và VeRA gây suy giảm PDR gián tiếp qua
overhead điều khiển.

\paragraph{End-to-End Delay (E2ED)}
E2ED là thời gian trung bình từ khi một gói dữ liệu được phát cho đến
khi đến đích thành công:
\begin{equation}
    \text{E2ED} = \frac{1}{N_{\text{received}}}
    \sum_{i=1}^{N_{\text{received}}} (t_{\text{receive},i} -
    t_{\text{send},i})
    \label{eq:e2ed}
\end{equation}
E2ED tăng khi mạng bị tắc nghẽn (DIS Flooding) hoặc khi gói phải truyền
qua đường dài hơn do tô-pô bị méo (Decreased Rank).

\paragraph{Control Overhead (CO)}
CO đo số lượng bản tin điều khiển RPL (DIO, DIS, DAO) phát sinh trong
một khoảng thời gian quan sát. Tấn công DIS Flooding và VeRA trực tiếp
làm tăng CO lên nhiều lần so với baseline.

\paragraph{Power Consumption}
Tiêu thụ năng lượng trung bình của mỗi nút (mW hoặc mJ), đo bằng công
cụ Powertrace trong Contiki-NG. Năng lượng tiêu thụ tăng khi nút phải
xử lý nhiều bản tin điều khiển hơn hoặc phải tái cấu trúc tô-pô thường
xuyên.

\subsubsection{Chỉ số đánh giá IDS}

\paragraph{True Positive Rate (TPR) -- Detection Rate}
Tỉ lệ các cuộc tấn công thực sự được IDS phát hiện đúng:
\begin{equation}
    \text{TPR} = \frac{TP}{TP + FN} \times 100\%
    \label{eq:tpr}
\end{equation}
trong đó $TP$ là số trường hợp tấn công phát hiện đúng, $FN$ là số
trường hợp tấn công bị bỏ sót.

\paragraph{False Positive Rate (FPR)}
Tỉ lệ các trường hợp bình thường bị IDS cảnh báo nhầm là tấn công:
\begin{equation}
    \text{FPR} = \frac{FP}{FP + TN} \times 100\%
    \label{eq:fpr}
\end{equation}
trong đó $FP$ là số cảnh báo nhầm, $TN$ là số trường hợp bình thường
được phân loại đúng. FPR cao làm tăng overhead và gây phiền nhiễu vận
hành.

\paragraph{Ma trận nhầm lẫn (Confusion Matrix)}
Tổng hợp kết quả phân loại của IDS qua bốn giá trị TP, TN, FP, FN như
Bảng~\ref{tab:confusion_matrix}.

\begin{table}[H]
\centering
\caption{Ma trận nhầm lẫn của IDS}
\label{tab:confusion_matrix}
\begin{tabularx}{0.7\textwidth}{C{1} C{1.2} C{1.2}}
    \toprule
    & \textbf{Thực tế: Tấn công} & \textbf{Thực tế: Bình thường} \\
    \midrule
    \textbf{Dự đoán: Tấn công}   & TP & FP \\
    \textbf{Dự đoán: Bình thường} & FN & TN \\
    \bottomrule
\end{tabularx}
\end{table}

% ============================================================
\subsection{Môi trường mô phỏng: Contiki-NG và Cooja}
\label{subsec:simulation_env}
% ============================================================

\subsubsection{Contiki-NG}

\textbf{Contiki-NG} là hệ điều hành mã nguồn mở, nhẹ nhàng dành riêng
cho các thiết bị IoT nhúng hạn chế tài nguyên. Contiki-NG là phiên bản
kế thừa và được tái cấu trúc hoàn toàn từ Contiki OS, với các cải tiến
về độ ổn định, bảo mật và khả năng kiểm thử. Một số đặc điểm nổi bật:

\begin{itemize}
    \item Hỗ trợ đầy đủ ngăn xếp IPv6/6LoWPAN và RPL theo RFC~6550.
    \item Lập lịch phân tán dựa trên protothreads, tiêu thụ rất ít bộ
          nhớ (vài KB RAM).
    \item Công cụ \textbf{Powertrace} cho phép đo tiêu thụ năng lượng
          chi tiết theo từng module (CPU, radio Tx, radio Rx...).
    \item Tích hợp sẵn Cooja Simulator và hỗ trợ cross-compile cho
          nhiều phần cứng (Tmote Sky, Zolertia RE-Mote, nRF52840...).
\end{itemize}

\subsubsection{Cooja Simulator}

\textbf{Cooja} là bộ mô phỏng mạng cảm biến tích hợp trong Contiki-NG,
cho phép mô phỏng hàng chục đến hàng trăm nút trên một máy tính thông
thường \cite{overview_rpl_framework}. Đặc điểm:

\begin{itemize}
    \item \textbf{Mô phỏng mức bytecode:} Chạy trực tiếp bytecode của
          Contiki-NG, đảm bảo hành vi hoàn toàn giống phần cứng thực.
    \item \textbf{Mô hình radio linh hoạt:} Hỗ trợ Unit Disk Graph
          Medium (UDGM -- lý tưởng) và UDGM Distance Loss (có suy hao
          theo khoảng cách). Đồ án sử dụng UDGM Distance Loss để mô
          phỏng thực tế hơn.
    \item \textbf{Scripting:} Hỗ trợ script JavaScript để tự động hóa
          thực nghiệm và thu thập kết quả.
    \item \textbf{Mote Output:} Ghi log chi tiết từng nút, bao gồm
          bản tin RPL gửi/nhận, thay đổi parent, giá trị Rank và cảnh
          báo IDS.
\end{itemize}

\subsubsection{Cấu hình mô phỏng chuẩn}

Bảng~\ref{tab:sim_config} tổng hợp các tham số mô phỏng chung được sử
dụng nhất quán trong tất cả các kịch bản kiểm thử của đồ án.

\begin{table}[H]
\centering
\caption{Tham số cấu hình mô phỏng Cooja}
\label{tab:sim_config}
\begin{tabularx}{\textwidth}{lX}
    \toprule
    \textbf{Tham số} & \textbf{Giá trị} \\
    \midrule
    Công cụ mô phỏng      & Contiki-NG + Cooja Simulator \\
    Mô hình radio          & UDGM: Distance Loss \\
    Số nút                 & 20 nút (1 Border Router + 19 mote) \\
    Phân bố không gian     & Ngẫu nhiên trong vùng 200m $\times$ 200m \\
    Objective Function     & MRHOF (ETX metric) \\
    Tần suất gửi dữ liệu  & 1 gói / 10 giây \\
    Thời gian mô phỏng     & 300 giây / kịch bản \\
    Tỉ lệ nút tấn công     & 0\%, 10\%, 25\%, 50\% \\
    Kịch bản đánh giá      & Baseline / Under Attack / With IDS \\
    \bottomrule
\end{tabularx}
\end{table}

% ============================================================
\subsection{Tóm tắt chương}
\label{subsec:ch2_summary}
% ============================================================

Chương này đã trình bày các nền tảng lý thuyết cần thiết cho đồ án, bao
gồm: (i)~nguyên lý hoạt động của RPL với cấu trúc DODAG, bốn loại bản
tin điều khiển, cơ chế Trickle Timer và hai Objective Function OF0/MRHOF;
(ii)~phân loại tấn công RPL theo ba nhóm Resources, Topology và Traffic;
(iii)~các phương pháp và kiến trúc IDS trong IoT/LLN; (iv)~bộ chỉ số
đánh giá hiệu năng PDR, E2ED, CO, Power, TPR, FPR; và (v)~môi trường
mô phỏng Contiki-NG/Cooja.

Các kiến thức này tạo nền tảng để phân tích chi tiết 4 dạng tấn công RPL
được trình bày trong Chương~3 và thiết kế IDS ở Chương~4.